\documentclass{article}

\title{Cross-Cultural Social Acceptability Classification via Cultural Context, Rating-Based Prompting, and Norm Retrieval}
\author{Anonymous ACL submission}
\date{}

\begin{document}

\maketitle

\begin{abstract}
Social acceptability judgments are not universal: the same behavior can be acceptable, unacceptable, or context-dependent depending on cultural norms, relationship roles, and situational details. We study whether language models can adapt their social judgments when given explicit cultural context. We construct a 120-example cross-cultural social acceptability classification dataset with three labels: acceptable, not acceptable, and context-dependent. Each example includes a scenario, a culture or social-context label, and a short cultural-context explanation. We evaluate GPT-based prompting strategies including scenario-only classification, cultural-context prompting, scalar acceptability rating, dataset-derived norm retrieval, and Social-Chem-101 retrieval with explicit intermediate norm findings. Our results show that cultural context is the largest driver of improvement. GPT-4o mini improves from 57.50 percent accuracy in the scenario-only baseline to 83.33 percent when cultural context is included in a rating-based prompt. The strongest GPT-4o mini setup combines cultural context, dataset-derived retrieval, and a 1--10 rating formulation, achieving 94.17 percent accuracy. A GPT-4o plus Social-Chem-101 RAG experiment achieves 89.17 percent accuracy while producing inspectable intermediate norm claims.
\end{abstract}

\section{Introduction}

Social acceptability classification asks whether a behavior, statement, or interaction is socially appropriate in a given context. This task is difficult because social norms are culturally situated. For example, calling a boss by their first name may be ordinary in an informal U.S. workplace but inappropriate in a setting with stronger hierarchy or formality expectations.

Large language models often produce fluent social judgments, but fluency does not guarantee cultural adaptability. A model may apply a generic or majority-culture norm even when the scenario provides a different cultural context. In this project, we ask whether models can adapt social acceptability judgments when the same behavior is paired with different cultural norms.

Compared with our earlier project plan, we remove DeBERTa from the main experimental story. Although DeBERTa was useful as a preliminary zero-shot NLI baseline, we did not train or fine-tune it on our task. Therefore, its results are less informative for our final research question than the GPT, RAG, and open-source LLM experiments.

\section{Related Work}

Prior work has shown that language models and benchmarks often encode social assumptions from their data sources. NLPositionality discusses how datasets reflect the positionality of their creators and annotators. GeoMLAMA shows that commonsense knowledge varies across geographic and cultural contexts. Recent cultural and safety evaluation benchmarks, including SALAD-Bench and NORMAD, similarly highlight that model behavior may not generalize uniformly across social settings.

Our project also draws on retrieval-augmented generation, which improves model behavior by providing task-relevant external knowledge at inference time. Instead of retrieving factual passages, we retrieve social norms and rules of thumb. We also use Social-Chem-101, a large dataset of social and moral rules of thumb, as a broad norm source for one RAG experiment.

\section{Dataset}

We constructed a dataset of 120 manually curated social acceptability examples. Each example contains a scenario or interaction, a culture or social-context label, a cultural-context explanation, and one gold label from acceptable, not acceptable, or context-dependent.

The dataset covers domains such as family expectations, privacy and consent, workplace hierarchy, communication style, politeness, friendship, fairness, authority, care and harm, and commitment. Many scenarios are intentionally paired across cultural contexts so that the same behavior may receive different labels depending on the norm description.

\begin{center}
\begin{tabular}{l r}
\hline
Label & Count \\
\hline
acceptable & 23 \\
context-dependent & 47 \\
not acceptable & 50 \\
\hline
Total & 120 \\
\hline
\end{tabular}
\end{center}

The context-dependent label is important because many social judgments depend on missing or variable factors, such as relationship closeness, consent, public versus private setting, hierarchy, severity, frequency, or prior agreement.

\section{Methods}

\subsection{Scenario-Only Baseline}

The first baseline gives the model only the scenario text and asks it to classify the behavior into one of the three labels. This tests whether the model can solve the task using generic social knowledge alone.

\subsection{Cultural-Context Prompting}

The next setup appends the dataset's culture label and cultural-context explanation to the prompt. This tests whether explicit cultural guidance improves classification. For example, instead of only seeing ``I called my boss by their first name,'' the model may also see a context such as ``Japan, formal-hierarchy norm'' with an explanation that formal address and seniority are important.

\subsection{Rating-Based Prompting}

Direct three-way classification can make models overcommit to a label too early. We therefore test a 1--10 scalar acceptability prompt. The model first assigns a score, where low scores indicate clear social unacceptability, high scores indicate clear acceptability, and middle scores indicate ambiguity or missing context. We then map scores to labels: 1--3 is not acceptable, 4--7 is context-dependent, and 8--10 is acceptable.

\subsection{Dataset-Derived RAG}

We construct controlled RAG databases from the dataset's cultural-context fields. These databases exclude gold labels and prior model predictions to avoid answer leakage. The retrieved chunks summarize cultural norms, scenario patterns, and conditions under which a behavior is acceptable, unacceptable, or context-dependent.

\subsection{Social-Chem-101 RAG with Intermediate Findings}

We also build a large RAG database from Social-Chem-101. We filter out low-quality rules of thumb, aggregate repeated annotations by rule ID, and preserve fields such as the rule of thumb, action, moral foundation, agreement, cultural pressure, legality, and inferred project-category hints. The full corpus contains 285,514 chunks, and the project-matched subset contains 265,743 chunks.

Unlike the dataset-derived RAG database, Social-Chem-101 is not culture-specific. We therefore treat it as broad social and moral norm evidence, while the dataset's cultural-context explanation remains the primary source for culture-specific judgment.

\section{Experiments and Results}

We report the most meaningful experiments rather than every ablation. The early DeBERTa baseline is omitted from the main table because it was not trained or fine-tuned for this task.

\begin{center}
\begin{tabular}{l l r}
\hline
Experiment & Method & Accuracy \\
\hline
Scenario-only baseline & GPT-4o mini, scenario only & 57.50 percent \\
Context baseline & GPT-4o mini, scenario plus context & 79.16 percent \\
Context plus rating & GPT-4o mini, context plus 1--10 rating & 83.33 percent \\
Context plus RAG v2 plus rating & GPT-4o mini, culture-aware RAG & 87.50 percent \\
Context plus RAG v1 plus rating & GPT-4o mini, dataset-derived RAG & 94.17 percent \\
Social-Chem RAG reasoning & GPT-4o, Social-Chem-101 RAG & 89.17 percent \\
Llama 3 8B & teammate results TBD & TBD \\
Qwen & teammate results TBD & TBD \\
\hline
\end{tabular}
\end{center}

The results show three main trends. First, scenario-only classification is weak because the model often applies a generic social judgment to culturally different cases. Second, adding explicit cultural context produces the largest single improvement. Third, rating-based prompting and retrieval can further improve accuracy, but only when the retrieved evidence is well matched to the scenario and culture.

The strongest result is the GPT-4o mini context plus dataset-derived RAG plus rating setup, with 94.17 percent accuracy. This suggests that a compact, task-aligned norm database can be more effective than a very large broad norm corpus. However, the Social-Chem-101 RAG reasoning experiment is still valuable because it exposes intermediate norm claims and makes the model's reasoning easier to inspect.

\section{Analysis}

\subsection{Why Cultural Context Helps}

The scenario-only baseline often fails because many examples are underdetermined without cultural information. Adding the cultural-context explanation gives the model a specific norm to condition on, reducing reliance on default assumptions.

\subsection{Why Rating-Based Prompting Helps}

The 1--10 rating prompt improves performance because it gives the model a way to represent borderline cases before mapping to a final label. Many social judgments are not naturally binary. A direct classifier must immediately choose among acceptable, not acceptable, and context-dependent, while the rating setup lets the model first express degree.

\subsection{Why Retrieval Helps Sometimes}

Retrieval helps when the retrieved norms match the scenario's key social variables. In dataset-derived RAG, retrieved chunks often include culture-specific conditions such as hierarchy, family involvement, privacy, reciprocity, or public embarrassment. However, retrieval can hurt when the retrieved evidence is too generic or points to a broad moral rule that does not capture the culture-specific nuance.

\subsection{Social-Chem-101 Reasoning Results}

The GPT-4o plus Social-Chem-101 RAG reasoning experiment achieved 89.17 percent accuracy. Its per-label performance was 82.61 percent for acceptable, 85.11 percent for context-dependent, and 96.00 percent for not acceptable. The model was strongest on clearly unacceptable cases, but it sometimes over-applied broad negative norms to cases labeled acceptable or context-dependent.

\section{Open-Source Model Experiments}

This section will be completed with teammate results.

\subsection{Llama 3 8B}

Add teammate setup here: model variant, prompt format, whether cultural context was included, whether examples or RAG were used, decoding settings, and final accuracy.

\subsection{Qwen}

Add teammate setup here: model variant, prompt format, whether cultural context was included, whether examples or RAG were used, decoding settings, and final accuracy.

\section{Conclusion}

Our experiments show that cross-cultural social acceptability classification is not simply a generic commonsense task. Models perform much better when they receive explicit cultural context, and additional gains are possible through rating-based prompting and task-aligned retrieval. The best GPT-4o mini configuration achieved 94.17 percent accuracy using scenario text, cultural context, dataset-derived RAG, and a 1--10 rating prompt. A GPT-4o Social-Chem-101 RAG experiment achieved 89.17 percent accuracy while producing intermediate norm findings, which improved interpretability even when accuracy was lower than the best task-specific RAG setup.

The main lesson is that culturally aware evaluation should test whether models adapt to the provided cultural norm, not merely whether they produce socially reasonable-sounding answers. Future work should expand the dataset to more cultures and domains, evaluate open-source models such as Llama 3 8B and Qwen, improve retrieval quality, and develop better methods for handling genuinely context-dependent cases.

\section*{References}

Forbes, Maxwell, Jena D. Hwang, Vered Shwartz, Maarten Sap, and Yejin Choi. 2020. Social Chemistry 101: Learning to Reason about Social and Moral Norms. EMNLP.

Lewis, Patrick, Ethan Perez, Aleksandra Piktus, Fabio Petroni, Vladimir Karpukhin, Naman Goyal, Heinrich Kuttler, Mike Lewis, Wen-tau Yih, Tim Rocktaschel, Sebastian Riedel, and Douwe Kiela. 2020. Retrieval-Augmented Generation for Knowledge-Intensive NLP Tasks.

Li, Lei et al. 2024. SALAD-Bench: A Hierarchical and Comprehensive Safety Benchmark for Large Language Models.

Rao et al. 2025. NORMAD: A Framework for Measuring the Cultural Adaptability of Large Language Models.

Santy, Sebastian et al. 2023. NLPositionality: Characterizing Design Biases of Datasets and Models.

Yin, Da et al. 2022. GeoMLAMA: Geo-Diverse Commonsense Probing on Multilingual Pre-Trained Language Models.

\end{document}
